	\documentclass[14pt]{extreport}
\usepackage{extsizes}
	\usepackage[frenchb]{babel}
	\usepackage[utf8]{inputenc}  
	\usepackage[T1]{fontenc}
	\usepackage{amssymb}
	\usepackage[mathscr]{euscript}
	\usepackage{stmaryrd}
	\usepackage{amsmath}
	\usepackage{tikz}
	\usepackage[all,cmtip]{xy}
	\usepackage{amsthm}
	\usepackage{varioref}
	\usepackage[ margin=1in]{geometry}
	\geometry{a4paper}
	\usepackage{lmodern}
	\usepackage{hyperref}
	\usepackage{array}
	\usepackage{float}
	\usepackage{easytable}
	 \usepackage{fancyhdr}\usepackage{longtable}
	 \usetikzlibrary{shapes.misc}
	 \newcommand\ang[1]{$#1{}^o$}
\newlength{\taillecellule}
\setlength{\taillecellule}{2cm}
\newcolumntype{C}{@{}>{\centering\arraybackslash}p{\taillecellule}@{}}

\renewcommand{\theenumi}{\alph{enumi})}
\usepackage{pstricks,multido}
\usepackage{arrayjob}
\usepackage{calc,xlop}
\tikzset{cross/.style={cross out, draw=black, minimum size=2*(#1-\pgflinewidth), inner sep=0pt, outer sep=0pt},
%default radius will be 1pt. 
cross/.default={1pt}}

	\pagestyle{fancy}
	\theoremstyle{plain}
	\fancyfoot[C]{\empty} 
	\fancyhead[L]{Bilan}
	\fancyhead[R]{Année de 5e}
	
	
	
	\title{Bilan année de 5e}
	\date{}
	\begin{document}



\subsection*{Exercice 1 : Calculs}  % 4 points
Recopiez et effectuez les calculs suivants : 

\begin{enumerate}
\item $2+3\times4\div2$
\item $5 - 2 + 3$
\item $ 3+2\times 4$
\item $ 5 + (-3) + 7 +  4 +  (-2) $
\item $ 3 - (-3) + 4 - 5 - (-2) + (-8)$
\item $(5-(-2) + 4) - (3- 2 + (-4))$
\end{enumerate}

\subsection*{Exercice 2 : Calculs aussi}

Effectuer : 
\begin{enumerate}
\item[a)] $ -5 + 4$
\item[b)] $ 11 - (- 5)$
\item[c)] $ - 7 - 3$
\item[d)] $ 2 + 5\times 4$
\item[e)] $ 5- 2 \times 2 +1$

\item[f)] $ -2 + 11 - (-2) + 7 - 3$
\item[g)] $ 3  - (-6 - 4)$
\item[h)] $ 12 - 5 - [3 - (8-9)] $
\item[i)] $ 	(2 \times 4- 13) + (6-  (-1))$
\item[j)] $ - 1 - 1 - (5 - 6)$
\end{enumerate}

\subsection*{Exercice 3 : Fractions}

\begin{enumerate}
\item Parmi la liste suivante, identifiez les fractions qui sont égales. Justifier les égalités \textbf{et} le fait que les groupes ainsi formés sont bien différents.
\[ \frac46;  \frac79; \frac34;  \frac{14}{21}; \frac12 \frac23; \frac{49}{63}; \frac{34}{68}\]
\item Trouver une fraction égale à $\frac34$ dont le dénominateur est $20$. 
\item Trouver une fraction égale à $\frac53$ dont le numérateur est $95$.
\item Trouver une fraction égale à $\frac69$ dont le numérateur est $10$.
\item Trouver une fraction égale à $\frac8{12}$ dont le dénominateur est $15$.
\item Trouver une fraction égale à $\frac59$ dont la somme du numérateur et du dénominateur est $70$. 
\item Trouver une fraction égale à $\frac73$ dont la différence entre le numérateur et le dénominateur est $48$. 
\item Trouver une fraction égale à $\frac{49}{37}$ dont le dénominateur soit compris entre $450$ et $490$.
\item Simplifier les fractions suivantes : 
\[ \frac5{10}; \frac9{15} ; \frac{35}{21} ; \frac{33}{88} ; \frac{52}{117} ; \frac{143}{187}\]

\item Simplifier les fractions 
\[ \frac{77}{121}; \phantom{meow}
\frac{156}{208};\phantom{meow}
\frac{225}{375};\phantom{meow}
\frac{125}{1~000};\phantom{meow}
\frac{130}{273}; \phantom{meow}
\frac{2~352}{5~376}; \phantom{meow}
\frac{30~752}{37~800}.\]
\item Réduire au même dénominateur les fractions : 
\[ \frac{17}{25}\text{  et  }\frac{19}{45};\phantom{meow}
\frac{5}{26}\text{  et  }\frac{7}{39};\phantom{meow}
\frac1{350}\text{  et  }\frac1{420};\phantom{meow}
\frac{13}{77}\text{  et  }\frac{9}{66}.\]
\item 
Réduisez les fractions suivantes au même dénominateur : $\frac23$, $\frac59$, $\frac56$, $\frac{11}{18}$. 

\item Trouver toutes les fractions égales à $\frac{60}{72}$ et dont les termes soient plus petits que ceux
de cette fraction. 

\item Calculer $\frac14+\frac12$ et $\frac12-\frac14$.
\item Calculer $\frac19+\frac13$ et $\frac13-\frac19$.
\item Calculer $\frac34 + \frac7{10}$ et $\frac34-\frac7{10}$.
\item Calculer $\frac5{24}+\frac38 - \frac3{10} + \frac1{20}$. 
 \item Calculer : 
 \[\frac13 + \frac23;  \phantom{meow}
 \frac34+\frac25; \phantom{meow}
 \frac23-\frac12;  \phantom{meow}
 \frac45-\frac23.\]
 \item Calculer : \[\left(4+\frac23\right) + \left(2+ \frac56\right);  \phantom{meow}
 \left(7 + \frac34\right) + \left(5 + \frac23\right)
 + \left(\frac13+\frac14\right).\]
 \item Calculer :\[ 17 - \left(\frac43 + \frac35\right) ;  \phantom{meow}
 \frac{15}4 + \left(\frac{17}9-\frac85\right); 
  \phantom{meow}
  \frac{13}{11} - \left(\frac79 - \frac23\right)\]
  \item Calculer : $\frac{31}{15} + \left(\frac{26}{35} - \frac{8}{21}\right)$.
\item Calculer : $\frac{91}{30} + \left(\frac{29}{42} - \frac{2}{35}\right)$.
\end{enumerate}




\subsection*{Exercice 4 - Conditions de constructibilité}

Parmi les données suivantes, lesquelles permettent de tracer un triangle ? Justifiez vos réponses. 

\begin{enumerate}
\item $AB = 3$ cm, $BC = 2$ cm, $AC = 6$ cm. 
\item $\widehat{ABC} =$ \ang{40}, $\widehat{BCA} =$ \ang{90}, $\widehat{CAB} =$ \ang{50}.
\item $AB = 5$ cm, $BC = 3$ cm, $AC = 7$ cm.
\item $\widehat{ABC} =$ \ang{30}, $\widehat{BCA} =$ \ang{50}, $\widehat{CAB} =$ \ang{90}.
\end{enumerate}

\subsection*{Exercice 5 - Tracé}

Parmi les consignes suivantes, est-il possible de construire un triangle ? Justifier vos réponses.
Quand la construction et possible, construire celui-ci à la règle, au compas et au rapporteur.
\begin{enumerate}
\item $AB= 3$ cm, $BC = 5$ cm, $AC = 1$ cm. 
\item $AB= 5$ cm, $BC = 6$ cm, et $AC = 8$ cm. 
\item $AB = 6$ cm, $AC = 13$ cm, $BC = 6$ cm.
\end{enumerate}
Construire le triangle ainsi choisi à la règle et au compas, puis construire : 
\begin{enumerate}
\item la médiatrice de $[AB]$
\item la hauteur issue de $C$
\item le symétrique du triangle $ABC$ par rapport à $(AB)$
\item le symétrique du triangle $ABC$ par rapport à $B$.
\end{enumerate} 
\subsection*{Exercice 6 - Angles}

\begin{enumerate}
\item Construire trois angles successivement adjacents : $\widehat{AOB}= 32^o$, $\widehat{BOC}= 72^o$, et $\widehat{COD}= 48^o$, puis les bissectrices $[OM)$, $[ON)$, $[OP)$ et $[OQ)$ des angles $\widehat{AOB}$, $\widehat{AOC}$, $\widehat{BOD}$ et $\widehat{COD}$. 
\item Calculer les angles $\widehat{MON}$ et $\widehat{POQ}$. Comparer
ces angles à l'angle $\widehat{BOC}$. 
\item Montrer que les angles $\widehat{MOQ}$ et $\widehat{NOP}$ ont la même bissectrice. 
\end{enumerate}

\subsection*{Exercice 7 - Propriété des angles}
	\begin{tikzpicture}[baseline]
    \tikzset{
      point/.style={thick, draw, cross out, inner sep=0pt, minimum width=5pt, minimum height=5pt,
      },
    }
    \clip (-1,-6.859954267158308) rectangle (9.893937016745921,3.298368333044762);
    	\draw[color={black}] (-49.337106354662005,-8.114802081529762)--(57.23104337140793,9.413170414574523);
	\draw[color={black}] (-43.01051039562473,25.49698162099757)--(50.800045009202414,-30.114681319318713);
	\draw[color={black}] (-32.54430446732648,-36.5053743530274)--(44.41902981165249,35.443898873669596);
	\draw[color={black}] (-49.10034384830023,-13.974751650033367)--(57.12663771366797,3.4970976845539177);
	\draw  [color={black},line width = 2,preaction={fill,color = {black},opacity = 0.2}] (0.9866000147758149,0.1631570570445111) -- (0,0) -- (0.8602101149709134,-0.5099395773234235) arc (-30.73:9.319999999999997:1) ;
	\draw [color={black}] (1.57,-0.3) node[anchor = center,scale=1, rotate = 0] {\ang{40}};
	\draw  [color={red},line width = 2,preaction={fill,color = {red},opacity = 0.2}] (4.841010222109848,-2.869763566876294) -- (3.980788327580091,-2.359843812402566) -- (4.711290183478223,-1.676933201590069) arc (42.82:-30.910000000000004:1) ;
	\draw [color={red}] (5.57,-2.19) node[anchor = center,scale=1, rotate = 0] {\ang{74}};
	\draw  [color={black},line width = 2,preaction={fill,color = {blue},opacity = 0.2}] (1.2237233194994683,-5.699013376196127) -- (0.23675925179005297,-5.859954267158308) -- (0.9672611725210718,-5.177043595736702) arc (42.82:9.009999999999998:1) ;
	\draw [color={black}] (1.67,-5.15) node[anchor = center,scale=1, rotate = 0] {\ang{34}}; 
	\draw [color={black}] (-0.5,0.5) node[anchor = center,scale=1, rotate = 0] {J};
	\draw [color={black}] (8.29,-4.62) node[anchor = center,scale=1, rotate = 0] {D};
	\draw [color={black}] (7.39,1.8) node[anchor = center,scale=1, rotate = 0] {W};
	\draw [color={black}] (3.98,-2.86) node[anchor = center,scale=1, rotate = 0] {U};
	\draw [color={black}] (0.74,-6.36) node[anchor = center,scale=1, rotate = 0] {I};

\end{tikzpicture}\\
\begin{enumerate}
\item  Déterminer la mesure de l'angle $\widehat{JUW}$. Justifiez.
\item  En déduire la mesure de l'angle $\widehat{UWJ}$.Justifiez.
\item  Déterminer si les droites $(JW)$ et $(ID)$ sont parallèles. Justifiez. 
\item En déduire la valeur de l'angle $\widehat{UDI}$.\end{enumerate}
\subsection*{Exercice 8 - Calcul littéral}
	Réduire les sommes suivantes : 
 \[ (x + 13) + (x - 9) - (x - 6); \phantom{meowmeow} 
 (a + 10) - (a - 5) + (a - 9). \]
 \[ x - (x - 7)  + (x - 9); \phantom{meowmeow} 
 a + (a - b + 7) - (a + b). \]
 \[ x + 13 - (x + y + 10) + (x - y); \phantom{meowmeow} 
 m - n - (7 - m) + (n - 9).\]
	
    \subsection*{Problèmes de réflexion plus avancés}
\begin{enumerate}
\item Quel nombre faut-il ajouter aux deux termes de la fraction $\frac7{13}$ pour qu'elle devienne égale à $\frac23$ ? 
 \item Quel nombre faut-il retrancher aux deux termes de la fraction $\frac{17}{23}$ pour obtenir une fraction égale à $\frac58$ ? 
 \item Quel nombre faut-il ajouter aux deux termes de la fraction $\frac58$ et retrancher aux deux termes de la fraction $\frac34$ pour obtenir deux fractions égales ? 
 \item Une personne dispose de deux heures pour effectuer une promenade. Elle part en tramway à la vitesse moyenne de $12$ km à l'heure et revient à pied, à la vitesse moyenne de $4$ km à l'heure. À quelle distance du point de départ devra-t-elle quitter le tramway ? 
 \item Un épicier achète un certain nombre de bouteilles de vin pour $225$ euros. En vendant la bouteille $1,65$ euros, il gagnerait autant que ce qu'il perdrait s'il la vendait $1,35$ euros. Quel est le nombre des bouteilles de vins ? 
 
 \item On considère un polygone régulier à $6$ côtés.
 \begin{enumerate}
 \item Réaliser une figure.
 \item Relier tous les sommets au centre $O$. 
 \item Montrer que tous les triangles ainsi formés ont leurs angles tous égaux à $60^o$. 
 \item En déduire que les angles d'un hexagone régulier mesurent $120$ degrés.
\end{enumerate}   

 \item On considère un polygone régulier à $5$ côtés.
 \begin{enumerate}
 \item Réaliser une figure.
 \item Relier tous les sommets au centre $O$. 
 \item Montrer que tous les triangles isocèles ainsi formés ont leurs angles à la base tous égaux à $54^o$. 
 \item En déduire que les angles d'un pentagone régulier mesurent $108^o$ degrés.
\end{enumerate}   
\item Utiliser les deux exercices précédents pour généraliser à un polygone régulier avec un nombre quelconque $N$ de côtés. On montrera ainsi que les angles d'une telle figure valent tous $\frac{N-2}N\times 180^o$.
\item On trace deux points $A$ et $B$ et une droite $(d)$ sur une feuille. Expliquer comment construire le cercle passant par $A$ et $B$ dont le centre appartient à $(d)$. Existe-t-il toujours ? Y en a-t-il toujours un seul ?
\item On considère un quadrilatère convexe $ABCD$ dans lequel la médiatrice de 
$[AB]$ est également la médiatrice de $[CD]$. Soit $O$ le point où cette médiatrice coupe la médiatrice de $[BC]$. \begin{enumerate}
\item Montrer que les quatre points $A$, $B$, $C$ et $D$ sont situés sur un même cercle de centre $O$. 
\item Démontrer que les médiatrices de $[AD]$, $[AC]$ et $[BD]$ passent également par $O$. 
\end{enumerate}
 \item Trouver trois nombres entiers consécutifs (consécutifs : qui se suivent) dont la somme est égale à $57$. 
 \item La somme de cinq nombres impairs consécutifs est 85. Quels sont ces nombres ? 
 \item Un père a 29 ans ; son fils a 5 ans ; dans combien d'années l'âge du père sera-t-il le triple de l'âge du fils ? 
 \item Un père a 41 ans ; ses trois enfants sont âgés de 7, 9 et 13 ans. Dans combien d'années l'âge du père sera-t-il égal à la somme des âges de ses enfants ? 
 \item Trouver un nombre dont la somme des quotients par 5, 7 et 9 soit égale à 429. 
 
\end{enumerate} 
 
\end{document}